% Part: normal-modal-logic
% Chapter: syntax-and-semantics
% Section: normal-modal-logics

\documentclass[../../../include/open-logic-section]{subfiles}

\begin{document}

\olfileid{nml}{syn}{nor}

\olsection{المنطقات الموجهية العادية}

يمكن النظر إلى المنطق الموجهي بوصفه أي مجموعة من~!!{formula}s تتمتع
بقدر مناسب من حسن السلوك. وثمة منطقات موجهية تحظى باهتمام خاص، لا لأن
لها تأويلات وتطبيقات مثيرة للاهتمام فحسب، بل أساسًا لأنها مفهومة جيدًا
من الناحية النظرية. وهي ما يسمى المنطقات الموجهية العادية.

\begin{defn}\ollabel{defn:modal-logic}
يكون~$\Log L$ \emph{منطقًا موجهيًا} إذا كان مجموعة من~!!{formula}s
اللغة الموجهية الأساسية بحيث
\begin{enumerate}
\item يحتوي جميع القضايا الصادقة منطقيًا،
\item ويكون مغلقًا تحت التعويض المنتظم،
\item ويكون مغلقًا تحت القياس الشرطي؛ أي كلما كان~$!A \in \Log L$ و
  $(!A \lif !B) \in \Log L$، كان أيضًا~$!B \in \Log L$.
\end{enumerate}
\end{defn}

\begin{ex}
من أمثلة مجموعات~!!{formula}s التي تحقق هذا التعريف، ومن ثم تُعد
منطقات موجهية، وإن لم تكن مثيرة جدًا للاهتمام:
\begin{enumerate}
\item مجموعة جميع القضايا الصادقة منطقيًا؛
\item ومجموعة جميع~!!{formula}s، أي المنطق الموجهي غير المتسق.
\end{enumerate}
\end{ex}

\begin{defn}\ollabel{defn:normal-modal-logic}
يكون المنطق الموجهي~$\Log L$ \emph{عاديًا} إذا وفقط إذا
\begin{enumerate}
\item كان يحتوي كلًا من
\begin{align}
\tag{K} \Box(p \lif q) & \lif (\Box p \lif \Box q) \\
\tag{Dual} \Diamond p & \liff \lnot \Box \lnot p
\end{align}
\item وكان مغلقًا تحت الإلزام؛ أي كلما كان~$!A \in \Log L$، كان أيضًا
  $\Box !A \in \Log L$.
\end{enumerate}
\end{defn}

سنصادف طريقتين رئيستين لتعريف المنطقات الموجهية. أولاهما برهانية
نظرية، بوصفها مجموعة~!!{formula}s التي يمكن اشتقاقها في أنظمة
!!{derivation} معينة. والثانية نموذجية نظرية، بوصفها مجموعات
!!{formula}s الصادقة في أصناف معينة من النماذج. وهذا يماثل الحال في
المنطق القضوي العادي، أو في منطق الرتبة الأولى أيضًا. فهنا كذلك يمكننا
تعيين مجموعة من~!!{formula}s بطريقتين مختلفتين، مثل تعيينها بوصفها
مجموعة القضايا الصادقة في جميع إسنادات قيم الصدق، وبوصفها مجموعة كل
!!{derivable}~!!{formula}s في نظام~!!{derivation} ما. وفي المنطقات
الموجهية العادية يمكن اتباع هاتين الطريقتين في التحديد باستخدام
النماذج العلائقية وأنظمة البرهان المسلّمية من نمط هيلبرت.
\end{document}
